\scanpage{56}
tenant compte des isom \inserted{can.} obtenus plus haut (de $A\otimes C$-modules)
\[
\mathfrak g\otimes_A C\simeq C_1,\qquad\check{\mathfrak g}^{*}\otimes_A C\simeq\prod_{\sigma\ne1}C_\sigma,
\]
\deleted{Il faut donc définir} \inserted{le théor. de Tate donne} des isom de $A\otimes_{\mathbb Z_p}C$-modules
\[
C_\sigma\longrightarrow C_\sigma,\quad\text{i.e.}\quad\lambda\longmapsto\lambda x_\sigma\quad(x_\sigma\in C)
\]
qui commutent à l'action de $\pi$ opérant sur la source par $*$, et sur le but par \deleted{\unsure{$g$}} $x\longmapsto g(x)$ si $\sigma\ne1$, par $x\longrightarrow\tau(\unsure{x|g})g(x)$ si $\sigma=1$ ($\tau$ étant la \unsure{contorsion} de Tate) Cela revient aux conditions
\[
\boxed{\left\{\begin{aligned}
g(x_\sigma)&=\sigma(\chi_\pi(g))x_\sigma&&\text{si }\sigma\ne1\\
g(x_1)&=\tau(g)^{-1}\chi_\pi(g)x_1&&\text{si }\sigma=1
\end{aligned}\right.\quad\Biggr\}\ \text{si }g\in\underset{\substack{\Vert\\\mathrm{Gal}(\overline K/K)}}{\pi}}
\]
\uline{NB} A priori, les $x_\sigma$ sont déterminés par la condition écrite modulo multiplication par un élém de $A^*\subset K^*\subset C^*$. Mais \unsure{le} (\unsure{sous participation seule|sans particularisation} du choix de $\overline K$) le choix de l'un de ces $x_\sigma$ \inserted{(disons, de $x_1$)} détermine canoniquement, par le théor. de Tate, tous les autres. \deleted{\unsure{De cette façon}}

Il n'existe pas de $x'_1\in C^*$ tel qu'on ait
\[
g(x'_1)=\chi_\pi(g)x'_1\qquad\text{pour }g\in\pi\ (\text{ni }=\text{ pour }J\subset\text{Inertie de }\pi)
\]
car faisant le produit de $x'_1$ et des $x_\sigma$ ($\sigma\ne1$), on trouverait $x\in C^*$ tel que
\[
g(x)=\underbrace{N_{A/\mathbb Z_p}(\chi_\pi(g))}_{=\tau(g)\text{ si }g\in\text{Inertie}}x,\qquad\text{impossible.}
\]
\mspage{55}
